% ====================================================================
%  Demonstration of the functionality of the 'mul-beamer' class
%
%  This file is included by talk.tex and lecture.tex via
%    % ====================================================================
%  Demonstration of the functionality of the 'mul-beamer' class
%
%  This file is included by talk.tex and lecture.tex via
%    % ====================================================================
%  Demonstration of the functionality of the 'mul-beamer' class
%
%  This file is included by talk.tex and lecture.tex via
%    % ====================================================================
%  Demonstration of the functionality of the 'mul-beamer' class
%
%  This file is included by talk.tex and lecture.tex via
%    \input{demo_functionality}
%  (the detailed documentation is in demo_documentation.tex)
%  Remove that line (and this file) when starting your own presentation.
% ====================================================================
%
% typesetting command names in the demo, e.g. \democmd{alert} -> \alert
\newcommand{\democmd}[1]{\texttt{\textbackslash #1}}
%
% agenda: table of contents in the layout of the MUL PowerPoint master
\agendaframe
%
%
\section{Overview}
\subsection{Class options}
% sketch of a footer style (scaled to half the text width): \demofooter{<option>}{<middle text>}{<slide number>};
% an empty <middle text> gives the footer without slogan (emptyFootings)
\newcommand{\demofooter}[3]{%
	\resizebox{0.48\textwidth}{!}{%
	\begin{tikzpicture}[x=1mm, y=1mm, every node/.style={inner sep=0pt, text=white, font=\fontsize{3.5}{4}\selectfont}]%
		\fill[digital_Turquoise] (0,0) rectangle (62,3.2);%
		\fill[digital_Turquoise] (0,1) -- (5,1) -- (0,6) -- cycle;%
		\ifblank{#2}{}{%
			\node[anchor=west] at (7,1.6) {\funneldisplay\bfseries\MULsloganText};%
			\node at (36,1.6) {#2};%
		}%
		\node[anchor=east] at (60,1.6) {#3};%
		\node[anchor=south west, text=digital_GraphiteBlack, font=\tiny\ttfamily] at (7,4) {#1};%
	\end{tikzpicture}}%
}%
\begin{frame}{Class options}
	\beamerdefaultoverlayspecification{}% all items at once
	\scriptsize
	\begin{columns}[T]
		\begin{column}{0.48\textwidth}
			\begin{description}
				\item[Footer] \texttt{defaultFootings}*, \texttt{fullFootings}, \texttt{lectureFootings}, \texttt{emptyFootings} (see below)
				\item[Slide numbers] \texttt{slideSingleNumber}*, \texttt{slideTotalNumber} (current/total)
				\item[Sidebar] \texttt{showSidebar}*, \texttt{hideSidebar}; \texttt{showothersubsections}*, \texttt{hideothersubsections}
				\item[Language] \texttt{DElocalization}*, \texttt{ENlocalization} (logo, claim, texts); \texttt{DElogo}, \texttt{ENlogo}, \texttt{DEslogan}, \texttt{ENslogan} (logo or claim only)
			\end{description}
		\end{column}
		\begin{column}{0.48\textwidth}
			\begin{description}
				\item[Title slide] \texttt{lightTitle}*, \texttt{darkTitle}
				\item[Text font] \texttt{classicFont}* (Latin Modern Sans), \texttt{brandFont} (Host Grotesk)
				\item[Results] \texttt{showResults}*, \texttt{hideResults} (\democmd{showResults})
				\item[beamer] all other options are passed to beamer, e.g.\ \texttt{trans}, \texttt{handout}; the aspect ratio is always 16:9
			\end{description}
		\end{column}
	\end{columns}
	\medskip
	* default
	\vfill
	\demofooter{defaultFootings}{F.~Author \textit{et al.} $|$ Department of Materials Science}{5}\hfill
	\demofooter{fullFootings}{F.~Author: Short title \quad Name of the Conference \quad 1 January 2026}{5/9}\par\medskip
	\demofooter{lectureFootings}{Name of the Course}{5}\hfill
	\demofooter{emptyFootings}{}{5}\par
	\smallskip
	{\tiny with \texttt{fullFootings} in lecture mode, the conference is replaced by the current section/subsection, e.g.\ \emph{1.2 Subsection}}
\end{frame}
%
\subsection{Commands}
\begin{frame}{Commands}
	\beamerdefaultoverlayspecification{}% all items at once
	\scriptsize
	\begin{columns}[T]
		\begin{column}{0.48\textwidth}
			\begin{description}
				\item[Title slide] \democmd{title[<short>]}, \democmd{author[<short>]} with \democmd{aff\{<n>\}}, \democmd{institute[<unit for footer>]}, \democmd{email}, \democmd{date}, \democmd{conference} (talk) or \democmd{lecture} (lecture mode), \democmd{additionsToTitlePage}
				\item[Settings] \democmd{setFrameLogo}, \democmd{setTitleLogo}, \democmd{setSlogan}, \democmd{setThanksText}, \democmd{setAgendaTitle}, \democmd{setDarkBackground}, \democmd{website}, \democmd{setQRcode}
				\item[Special slides] \democmd{agendaframe}, \democmd{personframe}, \democmd{finalframe} (between frames)
				\item[Lectures] \democmd{lectureBeamer} with \democmd{includeonlylecture} (beamer's original \democmd{lecture})
			\end{description}
		\end{column}
		\begin{column}{0.48\textwidth}
			\begin{description}
				\item[Content] \democmd{alert}, \democmd{Alert}, \democmd{example}, \democmd{funneldisplay}, \democmd{showResults}, colours \texttt{digital\_\dots}
				\item[Between frames] \democmd{literature}, \democmd{lastslide}
				\item[Placing] \texttt{overdraw} environment, \democmd{framegrid}
			\end{description}
		\end{column}
	\end{columns}
	\medskip
	All of them are described in \texttt{README.md} and in the backup slides at the end (\texttt{demo\_documentation.tex}); the following slides demonstrate them.
\end{frame}
%
%
\section{Basics}
\subsection{Text and highlighting}
\begin{frame}{Text and highlighting}
	\begin{itemize}[<1->]% all items at once (see "Lists and overlays")
		\item \alert{important content} is highlighted with \democmd{alert\{\dots\}}
		\item \Alert{less important content} is highlighted with \democmd{Alert\{\dots\}}
		\item {\funneldisplay The MUL display font} (Funnel Display) is switched on with \democmd{funneldisplay}; it is used for the frame titles
		\item the text is set in Latin Modern Sans; class option \texttt{brandFont} switches to Host Grotesk, the MUL text font
		\item mathematics works as usual: $a^2 + b^2 = \Alert{c^2}$, \quad $\sin\alpha = \frac{b}{c}$
	\end{itemize}
\end{frame}
%
\subsection{Blocks and columns}
\begin{frame}{Blocks and columns}
	\begin{columns}[T]
		\begin{column}{0.48\textwidth}
			\begin{block}{A block}
				Standard beamer \texttt{block} in the MUL colours.
			\end{block}
			\begin{alertblock}{An alert block}
				For warnings and important statements.
			\end{alertblock}
		\end{column}
		\begin{column}{0.48\textwidth}
			\begin{exampleblock}{An example block}
				For examples.
			\end{exampleblock}
			\begin{block}{}
				A block without title.
			\end{block}
		\end{column}
	\end{columns}
\end{frame}
%
\subsection{Colours}
% colour swatch with its name, e.g. \democolour{digital_Turquoise}{white}
\newcommand{\democolour}[2]{%
	\tikz[baseline=-0.6ex]\node[fill=#1, text=#2, font=\tiny\ttfamily, minimum width=26mm, minimum height=5mm, inner sep=1pt] {\detokenize{#1}};}
\begin{frame}{Colours}
	The brand colours are defined with the names of the MUL style guide (also used by the poster template):
	\medskip
	\begin{columns}[T]
		\begin{column}{0.32\textwidth}
			\textbf{primary}\par\smallskip
			\democolour{digital_Turquoise}{white}\par
			\democolour{digital_GraphiteBlack}{white}\par
			\democolour{digital_LightGreen}{digital_GraphiteBlack}\par
			\democolour{digital_TechGreen}{digital_GraphiteBlack}\par
			\democolour{digital_White}{digital_GraphiteBlack}
		\end{column}
		\begin{column}{0.32\textwidth}
			\textbf{combination} (with a primary colour)\par\smallskip
			\democolour{digital_BTLight}{white}\par
			\democolour{digital_BTDark}{white}\par
			\democolour{digital_TGLight}{white}\par
			\democolour{digital_TGDark}{white}\par
			\democolour{digital_GWLight}{digital_GraphiteBlack}\par
			\democolour{digital_GWDark}{digital_GraphiteBlack}
		\end{column}
		\begin{column}{0.32\textwidth}
			\textbf{secondary} (study areas only)\par\smallskip
			\democolour{digital_Violet}{white}\par
			\democolour{digital_Orange}{white}\par
			\democolour{digital_Blue}{white}
		\end{column}
	\end{columns}
	\medskip
	Use them as usual, e.g.\ \texttt{\textbackslash textcolor\{digital\_Turquoise\}\{\dots\}}.
\end{frame}
%
%
\section{Lists and overlays}
\subsection{Lists}
\begin{frame}{Lists}
	By default, list items are uncovered one by one (\texttt{<+->} is the class default):
	\begin{itemize}
		\item first level items have squares
		\begin{itemize}
			\item lower levels have dashes
		\end{itemize}
		\example examples are marked with \democmd{example} instead of \democmd{item}
		\item last item
	\end{itemize}
	\bigskip
	Use \texttt{\textbackslash begin\{itemize\}[<1->]} to show all items at once:
	\begin{itemize}[<1->]
		\item all
		\item at
		\item once
	\end{itemize}
\end{frame}
%
\subsection{Overlays}
\begin{frame}{Overlays}
	\begin{itemize}[<1->]
		\item \democmd{only<2>\{\dots\}}: content exists on slide 2 only (takes no space otherwise)
		\item \democmd{uncover<3->\{\dots\}}: space for the content is always reserved, the content is shown from slide 3 on
	\end{itemize}
	\bigskip
	\only<2>{\Alert{This text is shown with \democmd{only<2>}.}}
	\uncover<3->{\alert{This text is shown with \democmd{uncover<3->}.}}
\end{frame}
%
%
\section{Figures}
\subsection{Figures and literature}
\literature{%
	A. Author, B. Author, and C. Author, {\em J. Example Sci.} {\bf 1}, 1 (2026)\\
	D. Author and E. Author, {\em Example Lett.} {\bf 2}, 100 (2025)}
\begin{frame}{Figures and literature}
	\begin{columns}[T]
		\begin{column}{0.55\textwidth}
			\centering
			\input{figs/demo_plot}
		\end{column}
		\begin{column}{0.4\textwidth}
			\centering
			\input{figs/demo_lattice}
		\end{column}
	\end{columns}
	\bigskip
	\begin{itemize}[<1->]
		\item figures are kept in the folder \texttt{figs}, graphics (logos, background) in \texttt{graphics}
		\item images are included with \texttt{\textbackslash includegraphics\{figs/\dots\}}; these demo figures are TikZ drawings included with \texttt{\textbackslash input}
		\item \democmd{literature\{\dots\}} (outside a frame) prints references above the footer of this and all following slides; clear it with \democmd{literature\{\}}
	\end{itemize}
\end{frame}
\literature{}
%
\subsection{Placing objects}
\begin{frame}{Placing objects with \texttt{overdraw}}
	\begin{itemize}[<1->]
		\item the \texttt{overdraw} environment is a TikZ picture with absolute coordinates: one unit is 1/16 of the slide width, the origin is the bottom-left corner of the slide (right of the sidebar, if shown)
		\item \democmd{framegrid} inside \texttt{overdraw} draws the coordinate grid for finding positions -- remove it in the final version
		\item steps: grid (slide 2), grid with label and arrow placed with \texttt{overdraw} (slide 3), final result (slide 4)
	\end{itemize}
	\onslide<2-3>
	\begin{overdraw}
		\framegrid
	\end{overdraw}
	\onslide<3-4>
	\begin{overdraw}
		\node[draw=digital_Turquoise, very thick, rounded corners, fill=white, font=\small] (label) at (9,3) {placed at (9,3)};
		\draw[->, very thick, digital_Turquoise] (label) -- (5,4.5);
	\end{overdraw}
\end{frame}
%
%
\section{More tools}
\subsection{Hidden results}
\begin{frame}{Hidden results}
	\begin{itemize}[<1->]
		\item \democmd{showResults\{\dots\}} shows its content with the class option \texttt{showResults} (default), and leaves an empty space of the same size with \texttt{hideResults}
		\item independent of the beamer mode: in the presentation (default) as well as with the beamer class options \texttt{trans} or \texttt{handout} (one page per frame, without overlays, e.g.\ for printing), only \texttt{showResults}/\texttt{hideResults} decide
		\item handouts of exercises without results: \texttt{\textbackslash documentclass[trans, hideResults]\{mul-beamer\}}
	\end{itemize}
	\bigskip
	\begin{block}{Exercise}
		Solve $x^2 - 5x + 6 = 0$.\par\medskip
		Result: \showResults{$x_1 = 2$, \quad $x_2 = 3$}
	\end{block}
\end{frame}
%
\subsection{Frame logo}
% \setFrameLogo inside a group {...} changes the logo for the enclosed frames only;
% outside a group (or in the preamble) it changes it for all following frames
{%
\setFrameLogo{graphics/logo-dms.pdf}
\begin{frame}{Frame logo}
	\begin{itemize}[<1->]
		\item the logo in the top-right corner can be changed with \democmd{setFrameLogo\{graphics/\dots\}}
		\item here: the logo of the Department of Materials Science
	\end{itemize}
\end{frame}
}%
%
\subsection{Agenda slide}
\begin{frame}{Agenda slide}
	\begin{itemize}[<1->]
		\item \democmd{agendaframe} (between frames) shows the table of contents (sections only) in the \emph{Agenda} layout of the MUL PowerPoint master, see the second slide
		\item \democmd{agendaframe[<title>]} replaces the title (default: \emph{\MULagendaText}); \democmd{setAgendaTitle\{<title>\}} (in the preamble) changes the default
	\end{itemize}
\end{frame}
%
\subsection{Person slide}
\begin{frame}{Person slide}
	\begin{itemize}[<1->]
		\item \democmd{personframe[<photo>]\{<name>\}\{<role>\}\{<text>\}} (between frames) introduces a person, as in the \emph{Zur Person} layout of the MUL PowerPoint master
		\item the photo fills the left part of the slide (cropped as needed); without \texttt{<photo>}, a placeholder is shown (next slide)
		\item \texttt{<text>} can contain lists, e.g.\ a short CV
	\end{itemize}
\end{frame}
%
\personframe{Name Surname}{Function, Department of Materials Science}{%
	\begin{itemize}[<1->]
		\item since 2020: current position
		\item 2015--2020: previous position
		\item research interests: topic A, topic B
	\end{itemize}}
%
\subsection{Closing slide}
\begin{frame}{Closing slide}
	\begin{itemize}[<1->]
		\item \democmd{finalframe} (between frames) adds a closing slide in the dark style of the MUL slides -- see the end of \texttt{talk.tex}: ``Thank you for your attention!'' (localised), for talks the title and the authors, the e-mail address (\democmd{email}), the website (\democmd{website}, default unileoben.ac.at), an optional QR code (\democmd{setQRcode}), and the claim and author line in the footer
		\item \democmd{finalframe[\dots]} replaces the text, e.g.\ \texttt{\textbackslash finalframe[Questions?]}
		\item \democmd{setDarkBackground\{graphics/background\_dark.jpg\}} (in the preamble) uses the background of the official PowerPoint master instead of the CMS variant
		\item the class option \texttt{darkTitle} gives a title slide in the same dark style
	\end{itemize}
\end{frame}

%  (the detailed documentation is in demo_documentation.tex)
%  Remove that line (and this file) when starting your own presentation.
% ====================================================================
%
% typesetting command names in the demo, e.g. \democmd{alert} -> \alert
\newcommand{\democmd}[1]{\texttt{\textbackslash #1}}
%
% agenda: table of contents in the layout of the MUL PowerPoint master
\agendaframe
%
%
\section{Overview}
\subsection{Class options}
% sketch of a footer style (scaled to half the text width): \demofooter{<option>}{<middle text>}{<slide number>};
% an empty <middle text> gives the footer without slogan (emptyFootings)
\newcommand{\demofooter}[3]{%
	\resizebox{0.48\textwidth}{!}{%
	\begin{tikzpicture}[x=1mm, y=1mm, every node/.style={inner sep=0pt, text=white, font=\fontsize{3.5}{4}\selectfont}]%
		\fill[digital_Turquoise] (0,0) rectangle (62,3.2);%
		\fill[digital_Turquoise] (0,1) -- (5,1) -- (0,6) -- cycle;%
		\ifblank{#2}{}{%
			\node[anchor=west] at (7,1.6) {\funneldisplay\bfseries\MULsloganText};%
			\node at (36,1.6) {#2};%
		}%
		\node[anchor=east] at (60,1.6) {#3};%
		\node[anchor=south west, text=digital_GraphiteBlack, font=\tiny\ttfamily] at (7,4) {#1};%
	\end{tikzpicture}}%
}%
\begin{frame}{Class options}
	\beamerdefaultoverlayspecification{}% all items at once
	\scriptsize
	\begin{columns}[T]
		\begin{column}{0.48\textwidth}
			\begin{description}
				\item[Footer] \texttt{defaultFootings}*, \texttt{fullFootings}, \texttt{lectureFootings}, \texttt{emptyFootings} (see below)
				\item[Slide numbers] \texttt{slideSingleNumber}*, \texttt{slideTotalNumber} (current/total)
				\item[Sidebar] \texttt{showSidebar}*, \texttt{hideSidebar}; \texttt{showothersubsections}*, \texttt{hideothersubsections}
				\item[Language] \texttt{DElocalization}*, \texttt{ENlocalization} (logo, claim, texts); \texttt{DElogo}, \texttt{ENlogo}, \texttt{DEslogan}, \texttt{ENslogan} (logo or claim only)
			\end{description}
		\end{column}
		\begin{column}{0.48\textwidth}
			\begin{description}
				\item[Title slide] \texttt{lightTitle}*, \texttt{darkTitle}
				\item[Text font] \texttt{classicFont}* (Latin Modern Sans), \texttt{brandFont} (Host Grotesk)
				\item[Results] \texttt{showResults}*, \texttt{hideResults} (\democmd{showResults})
				\item[beamer] all other options are passed to beamer, e.g.\ \texttt{trans}, \texttt{handout}; the aspect ratio is always 16:9
			\end{description}
		\end{column}
	\end{columns}
	\medskip
	* default
	\vfill
	\demofooter{defaultFootings}{F.~Author \textit{et al.} $|$ Department of Materials Science}{5}\hfill
	\demofooter{fullFootings}{F.~Author: Short title \quad Name of the Conference \quad 1 January 2026}{5/9}\par\medskip
	\demofooter{lectureFootings}{Name of the Course}{5}\hfill
	\demofooter{emptyFootings}{}{5}\par
	\smallskip
	{\tiny with \texttt{fullFootings} in lecture mode, the conference is replaced by the current section/subsection, e.g.\ \emph{1.2 Subsection}}
\end{frame}
%
\subsection{Commands}
\begin{frame}{Commands}
	\beamerdefaultoverlayspecification{}% all items at once
	\scriptsize
	\begin{columns}[T]
		\begin{column}{0.48\textwidth}
			\begin{description}
				\item[Title slide] \democmd{title[<short>]}, \democmd{author[<short>]} with \democmd{aff\{<n>\}}, \democmd{institute[<unit for footer>]}, \democmd{email}, \democmd{date}, \democmd{conference} (talk) or \democmd{lecture} (lecture mode), \democmd{additionsToTitlePage}
				\item[Settings] \democmd{setFrameLogo}, \democmd{setTitleLogo}, \democmd{setSlogan}, \democmd{setThanksText}, \democmd{setAgendaTitle}, \democmd{setDarkBackground}, \democmd{website}, \democmd{setQRcode}
				\item[Special slides] \democmd{agendaframe}, \democmd{personframe}, \democmd{finalframe} (between frames)
				\item[Lectures] \democmd{lectureBeamer} with \democmd{includeonlylecture} (beamer's original \democmd{lecture})
			\end{description}
		\end{column}
		\begin{column}{0.48\textwidth}
			\begin{description}
				\item[Content] \democmd{alert}, \democmd{Alert}, \democmd{example}, \democmd{funneldisplay}, \democmd{showResults}, colours \texttt{digital\_\dots}
				\item[Between frames] \democmd{literature}, \democmd{lastslide}
				\item[Placing] \texttt{overdraw} environment, \democmd{framegrid}
			\end{description}
		\end{column}
	\end{columns}
	\medskip
	All of them are described in \texttt{README.md} and in the backup slides at the end (\texttt{demo\_documentation.tex}); the following slides demonstrate them.
\end{frame}
%
%
\section{Basics}
\subsection{Text and highlighting}
\begin{frame}{Text and highlighting}
	\begin{itemize}[<1->]% all items at once (see "Lists and overlays")
		\item \alert{important content} is highlighted with \democmd{alert\{\dots\}}
		\item \Alert{less important content} is highlighted with \democmd{Alert\{\dots\}}
		\item {\funneldisplay The MUL display font} (Funnel Display) is switched on with \democmd{funneldisplay}; it is used for the frame titles
		\item the text is set in Latin Modern Sans; class option \texttt{brandFont} switches to Host Grotesk, the MUL text font
		\item mathematics works as usual: $a^2 + b^2 = \Alert{c^2}$, \quad $\sin\alpha = \frac{b}{c}$
	\end{itemize}
\end{frame}
%
\subsection{Blocks and columns}
\begin{frame}{Blocks and columns}
	\begin{columns}[T]
		\begin{column}{0.48\textwidth}
			\begin{block}{A block}
				Standard beamer \texttt{block} in the MUL colours.
			\end{block}
			\begin{alertblock}{An alert block}
				For warnings and important statements.
			\end{alertblock}
		\end{column}
		\begin{column}{0.48\textwidth}
			\begin{exampleblock}{An example block}
				For examples.
			\end{exampleblock}
			\begin{block}{}
				A block without title.
			\end{block}
		\end{column}
	\end{columns}
\end{frame}
%
\subsection{Colours}
% colour swatch with its name, e.g. \democolour{digital_Turquoise}{white}
\newcommand{\democolour}[2]{%
	\tikz[baseline=-0.6ex]\node[fill=#1, text=#2, font=\tiny\ttfamily, minimum width=26mm, minimum height=5mm, inner sep=1pt] {\detokenize{#1}};}
\begin{frame}{Colours}
	The brand colours are defined with the names of the MUL style guide (also used by the poster template):
	\medskip
	\begin{columns}[T]
		\begin{column}{0.32\textwidth}
			\textbf{primary}\par\smallskip
			\democolour{digital_Turquoise}{white}\par
			\democolour{digital_GraphiteBlack}{white}\par
			\democolour{digital_LightGreen}{digital_GraphiteBlack}\par
			\democolour{digital_TechGreen}{digital_GraphiteBlack}\par
			\democolour{digital_White}{digital_GraphiteBlack}
		\end{column}
		\begin{column}{0.32\textwidth}
			\textbf{combination} (with a primary colour)\par\smallskip
			\democolour{digital_BTLight}{white}\par
			\democolour{digital_BTDark}{white}\par
			\democolour{digital_TGLight}{white}\par
			\democolour{digital_TGDark}{white}\par
			\democolour{digital_GWLight}{digital_GraphiteBlack}\par
			\democolour{digital_GWDark}{digital_GraphiteBlack}
		\end{column}
		\begin{column}{0.32\textwidth}
			\textbf{secondary} (study areas only)\par\smallskip
			\democolour{digital_Violet}{white}\par
			\democolour{digital_Orange}{white}\par
			\democolour{digital_Blue}{white}
		\end{column}
	\end{columns}
	\medskip
	Use them as usual, e.g.\ \texttt{\textbackslash textcolor\{digital\_Turquoise\}\{\dots\}}.
\end{frame}
%
%
\section{Lists and overlays}
\subsection{Lists}
\begin{frame}{Lists}
	By default, list items are uncovered one by one (\texttt{<+->} is the class default):
	\begin{itemize}
		\item first level items have squares
		\begin{itemize}
			\item lower levels have dashes
		\end{itemize}
		\example examples are marked with \democmd{example} instead of \democmd{item}
		\item last item
	\end{itemize}
	\bigskip
	Use \texttt{\textbackslash begin\{itemize\}[<1->]} to show all items at once:
	\begin{itemize}[<1->]
		\item all
		\item at
		\item once
	\end{itemize}
\end{frame}
%
\subsection{Overlays}
\begin{frame}{Overlays}
	\begin{itemize}[<1->]
		\item \democmd{only<2>\{\dots\}}: content exists on slide 2 only (takes no space otherwise)
		\item \democmd{uncover<3->\{\dots\}}: space for the content is always reserved, the content is shown from slide 3 on
	\end{itemize}
	\bigskip
	\only<2>{\Alert{This text is shown with \democmd{only<2>}.}}
	\uncover<3->{\alert{This text is shown with \democmd{uncover<3->}.}}
\end{frame}
%
%
\section{Figures}
\subsection{Figures and literature}
\literature{%
	A. Author, B. Author, and C. Author, {\em J. Example Sci.} {\bf 1}, 1 (2026)\\
	D. Author and E. Author, {\em Example Lett.} {\bf 2}, 100 (2025)}
\begin{frame}{Figures and literature}
	\begin{columns}[T]
		\begin{column}{0.55\textwidth}
			\centering
			% demonstration graphics: a simple plot drawn with TikZ
% (uses the plot colours defined by mul-beamer.cls)
\begin{tikzpicture}[x=0.9cm, y=0.9cm]
	% axes
	\draw[->] (0,0) -- (6.4,0) node[below left] {\scriptsize time};
	\draw[->] (0,0) -- (0,3.4) node[above right] {\scriptsize signal};
	\foreach \x in {1,...,6} {\draw (\x,0) -- (\x,-0.08);}
	% data
	\draw[very thick, plotblue, domain=0:6, samples=80, smooth] plot (\x, {1.6 + 1.2*sin(1.5*\x r)*exp(-0.25*\x)});
	\draw[very thick, plotorange, dashed, domain=0:6] plot (\x, {0.4 + 0.35*\x});
	\foreach \x/\y in {0.5/0.7, 1.5/1.0, 2.5/1.3, 3.5/1.5, 4.5/2.1, 5.5/2.3} {\fill[plotred] (\x,\y) circle (2pt);}
	% legend
	\node[anchor=north east, font=\scriptsize, align=left] at (6.4,3.4) {%
		\textcolor{plotblue}{\rule[0.5ex]{4mm}{1pt}} model A\\
		\textcolor{plotorange}{\rule[0.5ex]{4mm}{1pt}} model B\\
		\textcolor{plotred}{$\bullet$} measurement};
\end{tikzpicture}%

		\end{column}
		\begin{column}{0.4\textwidth}
			\centering
			% demonstration graphics: a schematic two-component lattice drawn with TikZ
\begin{tikzpicture}[x=0.55cm, y=0.55cm]
	\draw[gray!50] (-0.5,-0.5) rectangle (5.5,4.5);
	\foreach \i in {0,...,5} {
		\foreach \j in {0,...,4} {
			\pgfmathtruncatemacro{\site}{mod(\i+\j,2)}
			\ifnum\site=0
				\shade[ball color=plotblue] (\i,\j) circle (0.3);
			\else
				\shade[ball color=plotred] (\i,\j) circle (0.22);
			\fi
		}
	}
	\node[below, font=\scriptsize] at (2.5,-0.5) {%
		\textcolor{plotblue}{$\bullet$} species A\quad
		\textcolor{plotred}{$\bullet$} species B};
\end{tikzpicture}%

		\end{column}
	\end{columns}
	\bigskip
	\begin{itemize}[<1->]
		\item figures are kept in the folder \texttt{figs}, graphics (logos, background) in \texttt{graphics}
		\item images are included with \texttt{\textbackslash includegraphics\{figs/\dots\}}; these demo figures are TikZ drawings included with \texttt{\textbackslash input}
		\item \democmd{literature\{\dots\}} (outside a frame) prints references above the footer of this and all following slides; clear it with \democmd{literature\{\}}
	\end{itemize}
\end{frame}
\literature{}
%
\subsection{Placing objects}
\begin{frame}{Placing objects with \texttt{overdraw}}
	\begin{itemize}[<1->]
		\item the \texttt{overdraw} environment is a TikZ picture with absolute coordinates: one unit is 1/16 of the slide width, the origin is the bottom-left corner of the slide (right of the sidebar, if shown)
		\item \democmd{framegrid} inside \texttt{overdraw} draws the coordinate grid for finding positions -- remove it in the final version
		\item steps: grid (slide 2), grid with label and arrow placed with \texttt{overdraw} (slide 3), final result (slide 4)
	\end{itemize}
	\onslide<2-3>
	\begin{overdraw}
		\framegrid
	\end{overdraw}
	\onslide<3-4>
	\begin{overdraw}
		\node[draw=digital_Turquoise, very thick, rounded corners, fill=white, font=\small] (label) at (9,3) {placed at (9,3)};
		\draw[->, very thick, digital_Turquoise] (label) -- (5,4.5);
	\end{overdraw}
\end{frame}
%
%
\section{More tools}
\subsection{Hidden results}
\begin{frame}{Hidden results}
	\begin{itemize}[<1->]
		\item \democmd{showResults\{\dots\}} shows its content with the class option \texttt{showResults} (default), and leaves an empty space of the same size with \texttt{hideResults}
		\item independent of the beamer mode: in the presentation (default) as well as with the beamer class options \texttt{trans} or \texttt{handout} (one page per frame, without overlays, e.g.\ for printing), only \texttt{showResults}/\texttt{hideResults} decide
		\item handouts of exercises without results: \texttt{\textbackslash documentclass[trans, hideResults]\{mul-beamer\}}
	\end{itemize}
	\bigskip
	\begin{block}{Exercise}
		Solve $x^2 - 5x + 6 = 0$.\par\medskip
		Result: \showResults{$x_1 = 2$, \quad $x_2 = 3$}
	\end{block}
\end{frame}
%
\subsection{Frame logo}
% \setFrameLogo inside a group {...} changes the logo for the enclosed frames only;
% outside a group (or in the preamble) it changes it for all following frames
{%
\setFrameLogo{graphics/logo-dms.pdf}
\begin{frame}{Frame logo}
	\begin{itemize}[<1->]
		\item the logo in the top-right corner can be changed with \democmd{setFrameLogo\{graphics/\dots\}}
		\item here: the logo of the Department of Materials Science
	\end{itemize}
\end{frame}
}%
%
\subsection{Agenda slide}
\begin{frame}{Agenda slide}
	\begin{itemize}[<1->]
		\item \democmd{agendaframe} (between frames) shows the table of contents (sections only) in the \emph{Agenda} layout of the MUL PowerPoint master, see the second slide
		\item \democmd{agendaframe[<title>]} replaces the title (default: \emph{\MULagendaText}); \democmd{setAgendaTitle\{<title>\}} (in the preamble) changes the default
	\end{itemize}
\end{frame}
%
\subsection{Person slide}
\begin{frame}{Person slide}
	\begin{itemize}[<1->]
		\item \democmd{personframe[<photo>]\{<name>\}\{<role>\}\{<text>\}} (between frames) introduces a person, as in the \emph{Zur Person} layout of the MUL PowerPoint master
		\item the photo fills the left part of the slide (cropped as needed); without \texttt{<photo>}, a placeholder is shown (next slide)
		\item \texttt{<text>} can contain lists, e.g.\ a short CV
	\end{itemize}
\end{frame}
%
\personframe{Name Surname}{Function, Department of Materials Science}{%
	\begin{itemize}[<1->]
		\item since 2020: current position
		\item 2015--2020: previous position
		\item research interests: topic A, topic B
	\end{itemize}}
%
\subsection{Closing slide}
\begin{frame}{Closing slide}
	\begin{itemize}[<1->]
		\item \democmd{finalframe} (between frames) adds a closing slide in the dark style of the MUL slides -- see the end of \texttt{talk.tex}: ``Thank you for your attention!'' (localised), for talks the title and the authors, the e-mail address (\democmd{email}), the website (\democmd{website}, default unileoben.ac.at), an optional QR code (\democmd{setQRcode}), and the claim and author line in the footer
		\item \democmd{finalframe[\dots]} replaces the text, e.g.\ \texttt{\textbackslash finalframe[Questions?]}
		\item \democmd{setDarkBackground\{graphics/background\_dark.jpg\}} (in the preamble) uses the background of the official PowerPoint master instead of the CMS variant
		\item the class option \texttt{darkTitle} gives a title slide in the same dark style
	\end{itemize}
\end{frame}

%  (the detailed documentation is in demo_documentation.tex)
%  Remove that line (and this file) when starting your own presentation.
% ====================================================================
%
% typesetting command names in the demo, e.g. \democmd{alert} -> \alert
\newcommand{\democmd}[1]{\texttt{\textbackslash #1}}
%
% agenda: table of contents in the layout of the MUL PowerPoint master
\agendaframe
%
%
\section{Overview}
\subsection{Class options}
% sketch of a footer style (scaled to half the text width): \demofooter{<option>}{<middle text>}{<slide number>};
% an empty <middle text> gives the footer without slogan (emptyFootings)
\newcommand{\demofooter}[3]{%
	\resizebox{0.48\textwidth}{!}{%
	\begin{tikzpicture}[x=1mm, y=1mm, every node/.style={inner sep=0pt, text=white, font=\fontsize{3.5}{4}\selectfont}]%
		\fill[digital_Turquoise] (0,0) rectangle (62,3.2);%
		\fill[digital_Turquoise] (0,1) -- (5,1) -- (0,6) -- cycle;%
		\ifblank{#2}{}{%
			\node[anchor=west] at (7,1.6) {\funneldisplay\bfseries\MULsloganText};%
			\node at (36,1.6) {#2};%
		}%
		\node[anchor=east] at (60,1.6) {#3};%
		\node[anchor=south west, text=digital_GraphiteBlack, font=\tiny\ttfamily] at (7,4) {#1};%
	\end{tikzpicture}}%
}%
\begin{frame}{Class options}
	\beamerdefaultoverlayspecification{}% all items at once
	\scriptsize
	\begin{columns}[T]
		\begin{column}{0.48\textwidth}
			\begin{description}
				\item[Footer] \texttt{defaultFootings}*, \texttt{fullFootings}, \texttt{lectureFootings}, \texttt{emptyFootings} (see below)
				\item[Slide numbers] \texttt{slideSingleNumber}*, \texttt{slideTotalNumber} (current/total)
				\item[Sidebar] \texttt{showSidebar}*, \texttt{hideSidebar}; \texttt{showothersubsections}*, \texttt{hideothersubsections}
				\item[Language] \texttt{DElocalization}*, \texttt{ENlocalization} (logo, claim, texts); \texttt{DElogo}, \texttt{ENlogo}, \texttt{DEslogan}, \texttt{ENslogan} (logo or claim only)
			\end{description}
		\end{column}
		\begin{column}{0.48\textwidth}
			\begin{description}
				\item[Title slide] \texttt{lightTitle}*, \texttt{darkTitle}
				\item[Text font] \texttt{classicFont}* (Latin Modern Sans), \texttt{brandFont} (Host Grotesk)
				\item[Results] \texttt{showResults}*, \texttt{hideResults} (\democmd{showResults})
				\item[beamer] all other options are passed to beamer, e.g.\ \texttt{trans}, \texttt{handout}; the aspect ratio is always 16:9
			\end{description}
		\end{column}
	\end{columns}
	\medskip
	* default
	\vfill
	\demofooter{defaultFootings}{F.~Author \textit{et al.} $|$ Department of Materials Science}{5}\hfill
	\demofooter{fullFootings}{F.~Author: Short title \quad Name of the Conference \quad 1 January 2026}{5/9}\par\medskip
	\demofooter{lectureFootings}{Name of the Course}{5}\hfill
	\demofooter{emptyFootings}{}{5}\par
	\smallskip
	{\tiny with \texttt{fullFootings} in lecture mode, the conference is replaced by the current section/subsection, e.g.\ \emph{1.2 Subsection}}
\end{frame}
%
\subsection{Commands}
\begin{frame}{Commands}
	\beamerdefaultoverlayspecification{}% all items at once
	\scriptsize
	\begin{columns}[T]
		\begin{column}{0.48\textwidth}
			\begin{description}
				\item[Title slide] \democmd{title[<short>]}, \democmd{author[<short>]} with \democmd{aff\{<n>\}}, \democmd{institute[<unit for footer>]}, \democmd{email}, \democmd{date}, \democmd{conference} (talk) or \democmd{lecture} (lecture mode), \democmd{additionsToTitlePage}
				\item[Settings] \democmd{setFrameLogo}, \democmd{setTitleLogo}, \democmd{setSlogan}, \democmd{setThanksText}, \democmd{setAgendaTitle}, \democmd{setDarkBackground}, \democmd{website}, \democmd{setQRcode}
				\item[Special slides] \democmd{agendaframe}, \democmd{personframe}, \democmd{finalframe} (between frames)
				\item[Lectures] \democmd{lectureBeamer} with \democmd{includeonlylecture} (beamer's original \democmd{lecture})
			\end{description}
		\end{column}
		\begin{column}{0.48\textwidth}
			\begin{description}
				\item[Content] \democmd{alert}, \democmd{Alert}, \democmd{example}, \democmd{funneldisplay}, \democmd{showResults}, colours \texttt{digital\_\dots}
				\item[Between frames] \democmd{literature}, \democmd{lastslide}
				\item[Placing] \texttt{overdraw} environment, \democmd{framegrid}
			\end{description}
		\end{column}
	\end{columns}
	\medskip
	All of them are described in \texttt{README.md} and in the backup slides at the end (\texttt{demo\_documentation.tex}); the following slides demonstrate them.
\end{frame}
%
%
\section{Basics}
\subsection{Text and highlighting}
\begin{frame}{Text and highlighting}
	\begin{itemize}[<1->]% all items at once (see "Lists and overlays")
		\item \alert{important content} is highlighted with \democmd{alert\{\dots\}}
		\item \Alert{less important content} is highlighted with \democmd{Alert\{\dots\}}
		\item {\funneldisplay The MUL display font} (Funnel Display) is switched on with \democmd{funneldisplay}; it is used for the frame titles
		\item the text is set in Latin Modern Sans; class option \texttt{brandFont} switches to Host Grotesk, the MUL text font
		\item mathematics works as usual: $a^2 + b^2 = \Alert{c^2}$, \quad $\sin\alpha = \frac{b}{c}$
	\end{itemize}
\end{frame}
%
\subsection{Blocks and columns}
\begin{frame}{Blocks and columns}
	\begin{columns}[T]
		\begin{column}{0.48\textwidth}
			\begin{block}{A block}
				Standard beamer \texttt{block} in the MUL colours.
			\end{block}
			\begin{alertblock}{An alert block}
				For warnings and important statements.
			\end{alertblock}
		\end{column}
		\begin{column}{0.48\textwidth}
			\begin{exampleblock}{An example block}
				For examples.
			\end{exampleblock}
			\begin{block}{}
				A block without title.
			\end{block}
		\end{column}
	\end{columns}
\end{frame}
%
\subsection{Colours}
% colour swatch with its name, e.g. \democolour{digital_Turquoise}{white}
\newcommand{\democolour}[2]{%
	\tikz[baseline=-0.6ex]\node[fill=#1, text=#2, font=\tiny\ttfamily, minimum width=26mm, minimum height=5mm, inner sep=1pt] {\detokenize{#1}};}
\begin{frame}{Colours}
	The brand colours are defined with the names of the MUL style guide (also used by the poster template):
	\medskip
	\begin{columns}[T]
		\begin{column}{0.32\textwidth}
			\textbf{primary}\par\smallskip
			\democolour{digital_Turquoise}{white}\par
			\democolour{digital_GraphiteBlack}{white}\par
			\democolour{digital_LightGreen}{digital_GraphiteBlack}\par
			\democolour{digital_TechGreen}{digital_GraphiteBlack}\par
			\democolour{digital_White}{digital_GraphiteBlack}
		\end{column}
		\begin{column}{0.32\textwidth}
			\textbf{combination} (with a primary colour)\par\smallskip
			\democolour{digital_BTLight}{white}\par
			\democolour{digital_BTDark}{white}\par
			\democolour{digital_TGLight}{white}\par
			\democolour{digital_TGDark}{white}\par
			\democolour{digital_GWLight}{digital_GraphiteBlack}\par
			\democolour{digital_GWDark}{digital_GraphiteBlack}
		\end{column}
		\begin{column}{0.32\textwidth}
			\textbf{secondary} (study areas only)\par\smallskip
			\democolour{digital_Violet}{white}\par
			\democolour{digital_Orange}{white}\par
			\democolour{digital_Blue}{white}
		\end{column}
	\end{columns}
	\medskip
	Use them as usual, e.g.\ \texttt{\textbackslash textcolor\{digital\_Turquoise\}\{\dots\}}.
\end{frame}
%
%
\section{Lists and overlays}
\subsection{Lists}
\begin{frame}{Lists}
	By default, list items are uncovered one by one (\texttt{<+->} is the class default):
	\begin{itemize}
		\item first level items have squares
		\begin{itemize}
			\item lower levels have dashes
		\end{itemize}
		\example examples are marked with \democmd{example} instead of \democmd{item}
		\item last item
	\end{itemize}
	\bigskip
	Use \texttt{\textbackslash begin\{itemize\}[<1->]} to show all items at once:
	\begin{itemize}[<1->]
		\item all
		\item at
		\item once
	\end{itemize}
\end{frame}
%
\subsection{Overlays}
\begin{frame}{Overlays}
	\begin{itemize}[<1->]
		\item \democmd{only<2>\{\dots\}}: content exists on slide 2 only (takes no space otherwise)
		\item \democmd{uncover<3->\{\dots\}}: space for the content is always reserved, the content is shown from slide 3 on
	\end{itemize}
	\bigskip
	\only<2>{\Alert{This text is shown with \democmd{only<2>}.}}
	\uncover<3->{\alert{This text is shown with \democmd{uncover<3->}.}}
\end{frame}
%
%
\section{Figures}
\subsection{Figures and literature}
\literature{%
	A. Author, B. Author, and C. Author, {\em J. Example Sci.} {\bf 1}, 1 (2026)\\
	D. Author and E. Author, {\em Example Lett.} {\bf 2}, 100 (2025)}
\begin{frame}{Figures and literature}
	\begin{columns}[T]
		\begin{column}{0.55\textwidth}
			\centering
			% demonstration graphics: a simple plot drawn with TikZ
% (uses the plot colours defined by mul-beamer.cls)
\begin{tikzpicture}[x=0.9cm, y=0.9cm]
	% axes
	\draw[->] (0,0) -- (6.4,0) node[below left] {\scriptsize time};
	\draw[->] (0,0) -- (0,3.4) node[above right] {\scriptsize signal};
	\foreach \x in {1,...,6} {\draw (\x,0) -- (\x,-0.08);}
	% data
	\draw[very thick, plotblue, domain=0:6, samples=80, smooth] plot (\x, {1.6 + 1.2*sin(1.5*\x r)*exp(-0.25*\x)});
	\draw[very thick, plotorange, dashed, domain=0:6] plot (\x, {0.4 + 0.35*\x});
	\foreach \x/\y in {0.5/0.7, 1.5/1.0, 2.5/1.3, 3.5/1.5, 4.5/2.1, 5.5/2.3} {\fill[plotred] (\x,\y) circle (2pt);}
	% legend
	\node[anchor=north east, font=\scriptsize, align=left] at (6.4,3.4) {%
		\textcolor{plotblue}{\rule[0.5ex]{4mm}{1pt}} model A\\
		\textcolor{plotorange}{\rule[0.5ex]{4mm}{1pt}} model B\\
		\textcolor{plotred}{$\bullet$} measurement};
\end{tikzpicture}%

		\end{column}
		\begin{column}{0.4\textwidth}
			\centering
			% demonstration graphics: a schematic two-component lattice drawn with TikZ
\begin{tikzpicture}[x=0.55cm, y=0.55cm]
	\draw[gray!50] (-0.5,-0.5) rectangle (5.5,4.5);
	\foreach \i in {0,...,5} {
		\foreach \j in {0,...,4} {
			\pgfmathtruncatemacro{\site}{mod(\i+\j,2)}
			\ifnum\site=0
				\shade[ball color=plotblue] (\i,\j) circle (0.3);
			\else
				\shade[ball color=plotred] (\i,\j) circle (0.22);
			\fi
		}
	}
	\node[below, font=\scriptsize] at (2.5,-0.5) {%
		\textcolor{plotblue}{$\bullet$} species A\quad
		\textcolor{plotred}{$\bullet$} species B};
\end{tikzpicture}%

		\end{column}
	\end{columns}
	\bigskip
	\begin{itemize}[<1->]
		\item figures are kept in the folder \texttt{figs}, graphics (logos, background) in \texttt{graphics}
		\item images are included with \texttt{\textbackslash includegraphics\{figs/\dots\}}; these demo figures are TikZ drawings included with \texttt{\textbackslash input}
		\item \democmd{literature\{\dots\}} (outside a frame) prints references above the footer of this and all following slides; clear it with \democmd{literature\{\}}
	\end{itemize}
\end{frame}
\literature{}
%
\subsection{Placing objects}
\begin{frame}{Placing objects with \texttt{overdraw}}
	\begin{itemize}[<1->]
		\item the \texttt{overdraw} environment is a TikZ picture with absolute coordinates: one unit is 1/16 of the slide width, the origin is the bottom-left corner of the slide (right of the sidebar, if shown)
		\item \democmd{framegrid} inside \texttt{overdraw} draws the coordinate grid for finding positions -- remove it in the final version
		\item steps: grid (slide 2), grid with label and arrow placed with \texttt{overdraw} (slide 3), final result (slide 4)
	\end{itemize}
	\onslide<2-3>
	\begin{overdraw}
		\framegrid
	\end{overdraw}
	\onslide<3-4>
	\begin{overdraw}
		\node[draw=digital_Turquoise, very thick, rounded corners, fill=white, font=\small] (label) at (9,3) {placed at (9,3)};
		\draw[->, very thick, digital_Turquoise] (label) -- (5,4.5);
	\end{overdraw}
\end{frame}
%
%
\section{More tools}
\subsection{Hidden results}
\begin{frame}{Hidden results}
	\begin{itemize}[<1->]
		\item \democmd{showResults\{\dots\}} shows its content with the class option \texttt{showResults} (default), and leaves an empty space of the same size with \texttt{hideResults}
		\item independent of the beamer mode: in the presentation (default) as well as with the beamer class options \texttt{trans} or \texttt{handout} (one page per frame, without overlays, e.g.\ for printing), only \texttt{showResults}/\texttt{hideResults} decide
		\item handouts of exercises without results: \texttt{\textbackslash documentclass[trans, hideResults]\{mul-beamer\}}
	\end{itemize}
	\bigskip
	\begin{block}{Exercise}
		Solve $x^2 - 5x + 6 = 0$.\par\medskip
		Result: \showResults{$x_1 = 2$, \quad $x_2 = 3$}
	\end{block}
\end{frame}
%
\subsection{Frame logo}
% \setFrameLogo inside a group {...} changes the logo for the enclosed frames only;
% outside a group (or in the preamble) it changes it for all following frames
{%
\setFrameLogo{graphics/logo-dms.pdf}
\begin{frame}{Frame logo}
	\begin{itemize}[<1->]
		\item the logo in the top-right corner can be changed with \democmd{setFrameLogo\{graphics/\dots\}}
		\item here: the logo of the Department of Materials Science
	\end{itemize}
\end{frame}
}%
%
\subsection{Agenda slide}
\begin{frame}{Agenda slide}
	\begin{itemize}[<1->]
		\item \democmd{agendaframe} (between frames) shows the table of contents (sections only) in the \emph{Agenda} layout of the MUL PowerPoint master, see the second slide
		\item \democmd{agendaframe[<title>]} replaces the title (default: \emph{\MULagendaText}); \democmd{setAgendaTitle\{<title>\}} (in the preamble) changes the default
	\end{itemize}
\end{frame}
%
\subsection{Person slide}
\begin{frame}{Person slide}
	\begin{itemize}[<1->]
		\item \democmd{personframe[<photo>]\{<name>\}\{<role>\}\{<text>\}} (between frames) introduces a person, as in the \emph{Zur Person} layout of the MUL PowerPoint master
		\item the photo fills the left part of the slide (cropped as needed); without \texttt{<photo>}, a placeholder is shown (next slide)
		\item \texttt{<text>} can contain lists, e.g.\ a short CV
	\end{itemize}
\end{frame}
%
\personframe{Name Surname}{Function, Department of Materials Science}{%
	\begin{itemize}[<1->]
		\item since 2020: current position
		\item 2015--2020: previous position
		\item research interests: topic A, topic B
	\end{itemize}}
%
\subsection{Closing slide}
\begin{frame}{Closing slide}
	\begin{itemize}[<1->]
		\item \democmd{finalframe} (between frames) adds a closing slide in the dark style of the MUL slides -- see the end of \texttt{talk.tex}: ``Thank you for your attention!'' (localised), for talks the title and the authors, the e-mail address (\democmd{email}), the website (\democmd{website}, default unileoben.ac.at), an optional QR code (\democmd{setQRcode}), and the claim and author line in the footer
		\item \democmd{finalframe[\dots]} replaces the text, e.g.\ \texttt{\textbackslash finalframe[Questions?]}
		\item \democmd{setDarkBackground\{graphics/background\_dark.jpg\}} (in the preamble) uses the background of the official PowerPoint master instead of the CMS variant
		\item the class option \texttt{darkTitle} gives a title slide in the same dark style
	\end{itemize}
\end{frame}

%  (the detailed documentation is in demo_documentation.tex)
%  Remove that line (and this file) when starting your own presentation.
% ====================================================================
%
% typesetting command names in the demo, e.g. \democmd{alert} -> \alert
\newcommand{\democmd}[1]{\texttt{\textbackslash #1}}
%
% agenda: table of contents in the layout of the MUL PowerPoint master
\agendaframe
%
%
\section{Overview}
\subsection{Class options}
% sketch of a footer style (scaled to half the text width): \demofooter{<option>}{<middle text>}{<slide number>};
% an empty <middle text> gives the footer without slogan (emptyFootings)
\newcommand{\demofooter}[3]{%
	\resizebox{0.48\textwidth}{!}{%
	\begin{tikzpicture}[x=1mm, y=1mm, every node/.style={inner sep=0pt, text=white, font=\fontsize{3.5}{4}\selectfont}]%
		\fill[digital_Turquoise] (0,0) rectangle (62,3.2);%
		\fill[digital_Turquoise] (0,1) -- (5,1) -- (0,6) -- cycle;%
		\ifblank{#2}{}{%
			\node[anchor=west] at (7,1.6) {\funneldisplay\bfseries\MULsloganText};%
			\node at (36,1.6) {#2};%
		}%
		\node[anchor=east] at (60,1.6) {#3};%
		\node[anchor=south west, text=digital_GraphiteBlack, font=\tiny\ttfamily] at (7,4) {#1};%
	\end{tikzpicture}}%
}%
\begin{frame}{Class options}
	\beamerdefaultoverlayspecification{}% all items at once
	\scriptsize
	\begin{columns}[T]
		\begin{column}{0.48\textwidth}
			\begin{description}
				\item[Footer] \texttt{defaultFootings}*, \texttt{fullFootings}, \texttt{lectureFootings}, \texttt{emptyFootings} (see below)
				\item[Slide numbers] \texttt{slideSingleNumber}*, \texttt{slideTotalNumber} (current/total)
				\item[Sidebar] \texttt{showSidebar}*, \texttt{hideSidebar}; \texttt{showothersubsections}*, \texttt{hideothersubsections}
				\item[Language] \texttt{DElocalization}*, \texttt{ENlocalization} (logo, claim, texts); \texttt{DElogo}, \texttt{ENlogo}, \texttt{DEslogan}, \texttt{ENslogan} (logo or claim only)
			\end{description}
		\end{column}
		\begin{column}{0.48\textwidth}
			\begin{description}
				\item[Title slide] \texttt{lightTitle}*, \texttt{darkTitle}
				\item[Text font] \texttt{classicFont}* (Latin Modern Sans), \texttt{brandFont} (Host Grotesk)
				\item[Results] \texttt{showResults}*, \texttt{hideResults} (\democmd{showResults})
				\item[beamer] all other options are passed to beamer, e.g.\ \texttt{trans}, \texttt{handout}; the aspect ratio is always 16:9
			\end{description}
		\end{column}
	\end{columns}
	\medskip
	* default
	\vfill
	\demofooter{defaultFootings}{F.~Author \textit{et al.} $|$ Department of Materials Science}{5}\hfill
	\demofooter{fullFootings}{F.~Author: Short title \quad Name of the Conference \quad 1 January 2026}{5/9}\par\medskip
	\demofooter{lectureFootings}{Name of the Course}{5}\hfill
	\demofooter{emptyFootings}{}{5}\par
	\smallskip
	{\tiny with \texttt{fullFootings} in lecture mode, the conference is replaced by the current section/subsection, e.g.\ \emph{1.2 Subsection}}
\end{frame}
%
\subsection{Commands}
\begin{frame}{Commands}
	\beamerdefaultoverlayspecification{}% all items at once
	\scriptsize
	\begin{columns}[T]
		\begin{column}{0.48\textwidth}
			\begin{description}
				\item[Title slide] \democmd{title[<short>]}, \democmd{author[<short>]} with \democmd{aff\{<n>\}}, \democmd{institute[<unit for footer>]}, \democmd{email}, \democmd{date}, \democmd{conference} (talk) or \democmd{lecture} (lecture mode), \democmd{additionsToTitlePage}
				\item[Settings] \democmd{setFrameLogo}, \democmd{setTitleLogo}, \democmd{setSlogan}, \democmd{setThanksText}, \democmd{setAgendaTitle}, \democmd{setDarkBackground}, \democmd{website}, \democmd{setQRcode}
				\item[Special slides] \democmd{agendaframe}, \democmd{personframe}, \democmd{finalframe} (between frames)
				\item[Lectures] \democmd{lectureBeamer} with \democmd{includeonlylecture} (beamer's original \democmd{lecture})
			\end{description}
		\end{column}
		\begin{column}{0.48\textwidth}
			\begin{description}
				\item[Content] \democmd{alert}, \democmd{Alert}, \democmd{example}, \democmd{funneldisplay}, \democmd{showResults}, colours \texttt{digital\_\dots}
				\item[Between frames] \democmd{literature}, \democmd{lastslide}
				\item[Placing] \texttt{overdraw} environment, \democmd{framegrid}
			\end{description}
		\end{column}
	\end{columns}
	\medskip
	All of them are described in \texttt{README.md} and in the backup slides at the end (\texttt{demo\_documentation.tex}); the following slides demonstrate them.
\end{frame}
%
%
\section{Basics}
\subsection{Text and highlighting}
\begin{frame}{Text and highlighting}
	\begin{itemize}[<1->]% all items at once (see "Lists and overlays")
		\item \alert{important content} is highlighted with \democmd{alert\{\dots\}}
		\item \Alert{less important content} is highlighted with \democmd{Alert\{\dots\}}
		\item {\funneldisplay The MUL display font} (Funnel Display) is switched on with \democmd{funneldisplay}; it is used for the frame titles
		\item the text is set in Latin Modern Sans; class option \texttt{brandFont} switches to Host Grotesk, the MUL text font
		\item mathematics works as usual: $a^2 + b^2 = \Alert{c^2}$, \quad $\sin\alpha = \frac{b}{c}$
	\end{itemize}
\end{frame}
%
\subsection{Blocks and columns}
\begin{frame}{Blocks and columns}
	\begin{columns}[T]
		\begin{column}{0.48\textwidth}
			\begin{block}{A block}
				Standard beamer \texttt{block} in the MUL colours.
			\end{block}
			\begin{alertblock}{An alert block}
				For warnings and important statements.
			\end{alertblock}
		\end{column}
		\begin{column}{0.48\textwidth}
			\begin{exampleblock}{An example block}
				For examples.
			\end{exampleblock}
			\begin{block}{}
				A block without title.
			\end{block}
		\end{column}
	\end{columns}
\end{frame}
%
\subsection{Colours}
% colour swatch with its name, e.g. \democolour{digital_Turquoise}{white}
\newcommand{\democolour}[2]{%
	\tikz[baseline=-0.6ex]\node[fill=#1, text=#2, font=\tiny\ttfamily, minimum width=26mm, minimum height=5mm, inner sep=1pt] {\detokenize{#1}};}
\begin{frame}{Colours}
	The brand colours are defined with the names of the MUL style guide (also used by the poster template):
	\medskip
	\begin{columns}[T]
		\begin{column}{0.32\textwidth}
			\textbf{primary}\par\smallskip
			\democolour{digital_Turquoise}{white}\par
			\democolour{digital_GraphiteBlack}{white}\par
			\democolour{digital_LightGreen}{digital_GraphiteBlack}\par
			\democolour{digital_TechGreen}{digital_GraphiteBlack}\par
			\democolour{digital_White}{digital_GraphiteBlack}
		\end{column}
		\begin{column}{0.32\textwidth}
			\textbf{combination} (with a primary colour)\par\smallskip
			\democolour{digital_BTLight}{white}\par
			\democolour{digital_BTDark}{white}\par
			\democolour{digital_TGLight}{white}\par
			\democolour{digital_TGDark}{white}\par
			\democolour{digital_GWLight}{digital_GraphiteBlack}\par
			\democolour{digital_GWDark}{digital_GraphiteBlack}
		\end{column}
		\begin{column}{0.32\textwidth}
			\textbf{secondary} (study areas only)\par\smallskip
			\democolour{digital_Violet}{white}\par
			\democolour{digital_Orange}{white}\par
			\democolour{digital_Blue}{white}
		\end{column}
	\end{columns}
	\medskip
	Use them as usual, e.g.\ \texttt{\textbackslash textcolor\{digital\_Turquoise\}\{\dots\}}.
\end{frame}
%
%
\section{Lists and overlays}
\subsection{Lists}
\begin{frame}{Lists}
	By default, list items are uncovered one by one (\texttt{<+->} is the class default):
	\begin{itemize}
		\item first level items have squares
		\begin{itemize}
			\item lower levels have dashes
		\end{itemize}
		\example examples are marked with \democmd{example} instead of \democmd{item}
		\item last item
	\end{itemize}
	\bigskip
	Use \texttt{\textbackslash begin\{itemize\}[<1->]} to show all items at once:
	\begin{itemize}[<1->]
		\item all
		\item at
		\item once
	\end{itemize}
\end{frame}
%
\subsection{Overlays}
\begin{frame}{Overlays}
	\begin{itemize}[<1->]
		\item \democmd{only<2>\{\dots\}}: content exists on slide 2 only (takes no space otherwise)
		\item \democmd{uncover<3->\{\dots\}}: space for the content is always reserved, the content is shown from slide 3 on
	\end{itemize}
	\bigskip
	\only<2>{\Alert{This text is shown with \democmd{only<2>}.}}
	\uncover<3->{\alert{This text is shown with \democmd{uncover<3->}.}}
\end{frame}
%
%
\section{Figures}
\subsection{Figures and literature}
\literature{%
	A. Author, B. Author, and C. Author, {\em J. Example Sci.} {\bf 1}, 1 (2026)\\
	D. Author and E. Author, {\em Example Lett.} {\bf 2}, 100 (2025)}
\begin{frame}{Figures and literature}
	\begin{columns}[T]
		\begin{column}{0.55\textwidth}
			\centering
			% demonstration graphics: a simple plot drawn with TikZ
% (uses the plot colours defined by mul-beamer.cls)
\begin{tikzpicture}[x=0.9cm, y=0.9cm]
	% axes
	\draw[->] (0,0) -- (6.4,0) node[below left] {\scriptsize time};
	\draw[->] (0,0) -- (0,3.4) node[above right] {\scriptsize signal};
	\foreach \x in {1,...,6} {\draw (\x,0) -- (\x,-0.08);}
	% data
	\draw[very thick, plotblue, domain=0:6, samples=80, smooth] plot (\x, {1.6 + 1.2*sin(1.5*\x r)*exp(-0.25*\x)});
	\draw[very thick, plotorange, dashed, domain=0:6] plot (\x, {0.4 + 0.35*\x});
	\foreach \x/\y in {0.5/0.7, 1.5/1.0, 2.5/1.3, 3.5/1.5, 4.5/2.1, 5.5/2.3} {\fill[plotred] (\x,\y) circle (2pt);}
	% legend
	\node[anchor=north east, font=\scriptsize, align=left] at (6.4,3.4) {%
		\textcolor{plotblue}{\rule[0.5ex]{4mm}{1pt}} model A\\
		\textcolor{plotorange}{\rule[0.5ex]{4mm}{1pt}} model B\\
		\textcolor{plotred}{$\bullet$} measurement};
\end{tikzpicture}%

		\end{column}
		\begin{column}{0.4\textwidth}
			\centering
			% demonstration graphics: a schematic two-component lattice drawn with TikZ
\begin{tikzpicture}[x=0.55cm, y=0.55cm]
	\draw[gray!50] (-0.5,-0.5) rectangle (5.5,4.5);
	\foreach \i in {0,...,5} {
		\foreach \j in {0,...,4} {
			\pgfmathtruncatemacro{\site}{mod(\i+\j,2)}
			\ifnum\site=0
				\shade[ball color=plotblue] (\i,\j) circle (0.3);
			\else
				\shade[ball color=plotred] (\i,\j) circle (0.22);
			\fi
		}
	}
	\node[below, font=\scriptsize] at (2.5,-0.5) {%
		\textcolor{plotblue}{$\bullet$} species A\quad
		\textcolor{plotred}{$\bullet$} species B};
\end{tikzpicture}%

		\end{column}
	\end{columns}
	\bigskip
	\begin{itemize}[<1->]
		\item figures are kept in the folder \texttt{figs}, graphics (logos, background) in \texttt{graphics}
		\item images are included with \texttt{\textbackslash includegraphics\{figs/\dots\}}; these demo figures are TikZ drawings included with \texttt{\textbackslash input}
		\item \democmd{literature\{\dots\}} (outside a frame) prints references above the footer of this and all following slides; clear it with \democmd{literature\{\}}
	\end{itemize}
\end{frame}
\literature{}
%
\subsection{Placing objects}
\begin{frame}{Placing objects with \texttt{overdraw}}
	\begin{itemize}[<1->]
		\item the \texttt{overdraw} environment is a TikZ picture with absolute coordinates: one unit is 1/16 of the slide width, the origin is the bottom-left corner of the slide (right of the sidebar, if shown)
		\item \democmd{framegrid} inside \texttt{overdraw} draws the coordinate grid for finding positions -- remove it in the final version
		\item steps: grid (slide 2), grid with label and arrow placed with \texttt{overdraw} (slide 3), final result (slide 4)
	\end{itemize}
	\onslide<2-3>
	\begin{overdraw}
		\framegrid
	\end{overdraw}
	\onslide<3-4>
	\begin{overdraw}
		\node[draw=digital_Turquoise, very thick, rounded corners, fill=white, font=\small] (label) at (9,3) {placed at (9,3)};
		\draw[->, very thick, digital_Turquoise] (label) -- (5,4.5);
	\end{overdraw}
\end{frame}
%
%
\section{More tools}
\subsection{Hidden results}
\begin{frame}{Hidden results}
	\begin{itemize}[<1->]
		\item \democmd{showResults\{\dots\}} shows its content with the class option \texttt{showResults} (default), and leaves an empty space of the same size with \texttt{hideResults}
		\item independent of the beamer mode: in the presentation (default) as well as with the beamer class options \texttt{trans} or \texttt{handout} (one page per frame, without overlays, e.g.\ for printing), only \texttt{showResults}/\texttt{hideResults} decide
		\item handouts of exercises without results: \texttt{\textbackslash documentclass[trans, hideResults]\{mul-beamer\}}
	\end{itemize}
	\bigskip
	\begin{block}{Exercise}
		Solve $x^2 - 5x + 6 = 0$.\par\medskip
		Result: \showResults{$x_1 = 2$, \quad $x_2 = 3$}
	\end{block}
\end{frame}
%
\subsection{Frame logo}
% \setFrameLogo inside a group {...} changes the logo for the enclosed frames only;
% outside a group (or in the preamble) it changes it for all following frames
{%
\setFrameLogo{graphics/logo-dms.pdf}
\begin{frame}{Frame logo}
	\begin{itemize}[<1->]
		\item the logo in the top-right corner can be changed with \democmd{setFrameLogo\{graphics/\dots\}}
		\item here: the logo of the Department of Materials Science
	\end{itemize}
\end{frame}
}%
%
\subsection{Agenda slide}
\begin{frame}{Agenda slide}
	\begin{itemize}[<1->]
		\item \democmd{agendaframe} (between frames) shows the table of contents (sections only) in the \emph{Agenda} layout of the MUL PowerPoint master, see the second slide
		\item \democmd{agendaframe[<title>]} replaces the title (default: \emph{\MULagendaText}); \democmd{setAgendaTitle\{<title>\}} (in the preamble) changes the default
	\end{itemize}
\end{frame}
%
\subsection{Person slide}
\begin{frame}{Person slide}
	\begin{itemize}[<1->]
		\item \democmd{personframe[<photo>]\{<name>\}\{<role>\}\{<text>\}} (between frames) introduces a person, as in the \emph{Zur Person} layout of the MUL PowerPoint master
		\item the photo fills the left part of the slide (cropped as needed); without \texttt{<photo>}, a placeholder is shown (next slide)
		\item \texttt{<text>} can contain lists, e.g.\ a short CV
	\end{itemize}
\end{frame}
%
\personframe{Name Surname}{Function, Department of Materials Science}{%
	\begin{itemize}[<1->]
		\item since 2020: current position
		\item 2015--2020: previous position
		\item research interests: topic A, topic B
	\end{itemize}}
%
\subsection{Closing slide}
\begin{frame}{Closing slide}
	\begin{itemize}[<1->]
		\item \democmd{finalframe} (between frames) adds a closing slide in the dark style of the MUL slides -- see the end of \texttt{talk.tex}: ``Thank you for your attention!'' (localised), for talks the title and the authors, the e-mail address (\democmd{email}), the website (\democmd{website}, default unileoben.ac.at), an optional QR code (\democmd{setQRcode}), and the claim and author line in the footer
		\item \democmd{finalframe[\dots]} replaces the text, e.g.\ \texttt{\textbackslash finalframe[Questions?]}
		\item \democmd{setDarkBackground\{graphics/background\_dark.jpg\}} (in the preamble) uses the background of the official PowerPoint master instead of the CMS variant
		\item the class option \texttt{darkTitle} gives a title slide in the same dark style
	\end{itemize}
\end{frame}
